\section{Wnioski}

\subsection{Podsumowanie wyników}

Projekt zrealizował cel przewidywania długości pobytu w szpitalu. Przetestowano 3 architektury (SimpleMLP, TabNet, TabTransformer) z łącznie 62 konfiguracjami. Wszystkie modele osiągnęły R² powyżej 0.64.

Najlepszy model (TabNet): MAE = 1.26 dnia, R² = 0.67.

\subsection{Stosowalność rozwiązania}

Model może wspierać:
\begin{itemize}
    \item Planowanie wykorzystania łóżek szpitalnych
    \item Optymalizację zarządzania personelem
    \item Analizę kosztów i budżetowanie
\end{itemize}

Model powinien wspomagać, nie zastępować osąd medyczny.

\subsection{Ograniczenia}

\begin{itemize}
    \item Ograniczona liczba cech (22)
    \item Trudności z ekstremalnymi czasami pobytu
    \item 33\% niewytłumaczonej wariancji
    \item Brak walidacji na danych z innych szpitali
\end{itemize}

\subsection{Kierunki przyszłych badań}

\begin{itemize}
    \item Wzbogacenie danych o komplikacje i procedury medyczne
    \item Modelowanie niepewności (przedziały ufności)
    \item Walidacja w rzeczywistym środowisku szpitalnym
    \item Integracja z systemami HIS
\end{itemize}

\subsection{Wnioski końcowe}

\begin{enumerate}
    \item TabNet jest najlepszym modelem (MAE = 1.26 dnia, R² = 0.67)
    \item Grid search z 5-fold CV okazał się efektywną metodą optymalizacji
    \item Specjalizowane architektury przewyższają klasyczny MLP o 3-4\%
    \item Interpretowalność TabNet czyni go wartościowym w kontekście medycznym
    \item Model jest gotowy do wdrożenia jako narzędzie wspomagające
\end{enumerate}

Projekt wykazał, że sieci neuronowe mogą skutecznie wspierać planowanie zarządzania zasobami szpitalnymi.

